\chapter{Deployment Strategy and Distributed Execution}

\section{Introduction}
This chapter explains how \projectname{} is physically deployed and how its distributed execution is realized in practice. It clarifies the separation of services, the movement of data and control across the environment, and the operational mechanisms used to reproduce the deployment. The approach combines distributed-service reasoning, containerized packaging, dependency-aware startup logic, and remote execution paths \cite{docker_docs,apache_spark_docs,apache_hadoop_docs}.

\section{Control Plane Responsibilities}
The project adopts a deployment strategy based on a clear separation between a \textbf{local control plane} and a \textbf{distributed data plane}. The control plane is hosted on the local workstation, while the heavy data services are distributed across the three Linux virtual machines.

\begin{table}[!htbp]
    \centering
    \footnotesize
    \renewcommand{\arraystretch}{1.16}
    \arrayrulecolor{tablelineblue}
    \rowcolors{2}{tablebodyblue}{white}
    \begin{tabularx}{\textwidth}{|Y|Y|Y|}
        \hline
        \rowcolor{tableheaderblue}
        \textcolor{white}{\textbf{Control-plane service}} & \textcolor{white}{\textbf{Operational role}} & \textcolor{white}{\textbf{Reason for local placement}} \\
        \hline
        Airflow & orchestration of reset, setup, readiness, integrated raw loading/quality, and dbt transformation DAGs & avoids consuming VM resources for workflow control and keeps execution coordination centralized \\
        \hline
        Prometheus & collection of infrastructure and pipeline metrics & centralizes observability and simplifies metric management \\
        \hline
        Grafana & dashboard-based monitoring of infrastructure and pipeline state & provides one visual supervision point for the whole platform \\
        \hline
        Kafka UI & broker and topic inspection & keeps operational visibility outside the distributed data services themselves \\
        \hline
        Pipeline metrics exporter & exposure of pipeline-specific operational metrics & complements infrastructure metrics with workflow-level signals \\
        \hline
        Local TLS gateway & protected browser access to control-plane interfaces & centralizes HTTPS exposure for Airflow, Grafana, Kafka UI, and Prometheus \\
        \hline
        SSH submission path & remote Spark submission to the entry point on VM2 & avoids local-driver networking constraints and centralizes cluster entry \\
        \hline
    \end{tabularx}
    \caption{Main responsibilities of the local control plane}
    \label{tab:chapter6-control-plane}
\end{table}

This separation prevents the virtual machines from being loaded with orchestration and monitoring services that do not participate directly in distributed storage or processing. It also establishes a clear responsibility boundary: the control plane orchestrates, supervises, and initiates remote work, while the data plane transports, stores, processes, governs, and serves data.

\begin{figure}[H]
    \centering
    \begin{tikzpicture}[node distance=0.6cm and 0.6cm]
    \node[ado box=adograyfill, text width=10.6cm, minimum height=1.4cm] (cp)
    {\adoboxcontent{Local control plane}{orchestration + monitoring + supervision}};

    \node[ado box=adobluefill, text width=2.7cm, below left=1.0cm and -0.1cm of cp.south] (vm1)
    {\adoboxcontent{VM1}{broker 1\\DataNode 1\\Spark worker 1\\Trino worker 1}};
    \node[ado box=adotealfill, text width=3.2cm, below=1.0cm of cp.south] (vm2)
    {\adoboxcontent{VM2}{broker 2\\NameNode\\Hive Metastore\\Spark master\\Trino coordinator}};
    \node[ado box=adopurplefill, text width=2.7cm, below right=1.0cm and -0.1cm of cp.south] (vm3)
    {\adoboxcontent{VM3}{broker 3\\DataNode 3\\Spark worker 3\\Trino worker 2}};

    \node[ado group=adoborder, fit=(vm1)(vm2)(vm3), inner sep=0.35cm, label={[ado label]north:Distributed Linux data plane}] (plane) {};

    \draw[ado arrow] (cp.south) -- (plane.north);
    \draw[ado dashed arrow] (cp.south west) .. controls +(down:0.9cm) and +(up:0.6cm) .. (vm1.north);
    \draw[ado dashed arrow] (cp.south) -- (vm2.north);
    \draw[ado dashed arrow] (cp.south east) .. controls +(down:0.9cm) and +(up:0.6cm) .. (vm3.north);
\end{tikzpicture}

    \caption{Control-plane supervision of the distributed Linux data plane.}
    \label{fig:chapter6-control-vs-data-plane}
\end{figure}
\FloatBarrier

\section{VM-Level Service Separation}
The data plane is not deployed as three identical generic machines. Instead, the services are deliberately separated according to their coordination role or worker role.

\begin{table}[!htbp]
    \centering
    \footnotesize
    \renewcommand{\arraystretch}{1.16}
    \arrayrulecolor{tablelineblue}
    \rowcolors{2}{tablebodyblue}{white}
    \begin{tabularx}{\textwidth}{|L{1.4cm}|Y|Y|}
        \hline
        \rowcolor{tableheaderblue}
        \textcolor{white}{\textbf{Node}} & \textcolor{white}{\textbf{Main hosted services}} & \textcolor{white}{\textbf{Architectural logic}} \\
        \hline
        VM1 & Kafka broker 1, HDFS DataNode 1, Spark worker 1, Trino worker 1, cAdvisor & extends storage, compute, and query-worker capacity without carrying cluster leadership responsibilities \\
        \hline
        VM2 & Kafka broker 2, Kerberos KDC, NameNode, DataNode 2, Hive Metastore and PostgreSQL store, Spark master, Spark worker 2, Spark submit client, Spark Thrift Server, Spark History Server, Trino coordinator, Trino gateway, cAdvisor & acts as the leader node and integration point for coordination-heavy services \\
        \hline
        VM3 & Kafka broker 3, HDFS DataNode 3, Spark worker 3, Trino worker 2, cAdvisor & reinforces the cluster with additional worker capacity and distributed persistence \\
        \hline
    \end{tabularx}
    \caption{VM-level service allocation in the distributed deployment}
    \label{tab:chapter6-vm-separation}
\end{table}

This service distribution supports clear operational responsibility. VM2 concentrates services that require a central authority or coordination function, including Kerberos, the HDFS namespace, the shared catalog, Spark scheduling, application history, and distributed-query coordination. VM1 and VM3 provide horizontal extension through storage blocks, Spark execution resources, Kafka brokers, and Trino worker capacity.

Unlike a pseudo-distributed installation, the principal storage, processing, transport, and query responsibilities are materially distributed across separate machines. The deployment topology is therefore part of the implemented engineering contribution.

\begin{figure}[H]
    \centering
    \begin{tikzpicture}[x=1cm,y=1cm,font=\small]
  \node[ado group, minimum width=14.60cm, minimum height=7.10cm, anchor=north] (panel) at (0,0) {};
  \node[ado label, font=\bfseries\large] at ($(panel.north)+(0,-0.38)$)
    {VM service placement and node roles};

  \node[ado box={12.40cm}{adotealfill}, minimum height=2.10cm, inner sep=6pt] (vm2) at (0,-2.45)
    {\textbf{\small VM2 coordination node}\\[3pt]
     {\scriptsize
     \begin{tabular}{@{}p{5.55cm}p{5.55cm}@{}}
       \textbf{Security:} Kerberos KDC & \textbf{Catalog:} Hive Metastore + PostgreSQL \\
       \textbf{Transport:} Kafka broker 2 & \textbf{HDFS:} NameNode + DataNode 2 \\
       \textbf{Spark core:} master + worker 2 & \textbf{Spark access:} submit client + Thrift \\
       \textbf{Query:} Trino coordinator & \textbf{Spark evidence:} History Server \\
       \textbf{Operations:} cAdvisor & \textbf{Access:} Trino gateway
     \end{tabular}}};

  \node[ado box={5.75cm}{adobluefill}, minimum height=1.85cm, inner sep=6pt] (vm1) at (-3.35,-5.45)
    {\textbf{\small VM1 worker node}\\[3pt]
     \textbf{\scriptsize Kafka:} \scriptsize broker 1 \quad
     \textbf{\scriptsize HDFS:} \scriptsize DataNode 1\\[2pt]
     \textbf{\scriptsize Spark:} \scriptsize worker 1 \quad
     \textbf{\scriptsize Trino:} \scriptsize worker 1\\[2pt]
     \textbf{\scriptsize Operations:} \scriptsize cAdvisor};

  \node[ado box={5.75cm}{adopurplefill}, minimum height=1.85cm, inner sep=6pt] (vm3) at (3.35,-5.45)
    {\textbf{\small VM3 worker node}\\[3pt]
     \textbf{\scriptsize Kafka:} \scriptsize broker 3 \quad
     \textbf{\scriptsize HDFS:} \scriptsize DataNode 3\\[2pt]
     \textbf{\scriptsize Spark:} \scriptsize worker 3 \quad
     \textbf{\scriptsize Trino:} \scriptsize worker 2\\[2pt]
     \textbf{\scriptsize Operations:} \scriptsize cAdvisor};
\end{tikzpicture}

    \caption{Service placement and node roles across VM1, VM2, and VM3.}
    \label{fig:chapter6-vm-service-map}
\end{figure}
\FloatBarrier

\section{Data Flow Across Machines}
The implemented machine-to-machine flow follows an explicit sequence:
\begin{enumerate}[leftmargin=*,itemsep=2pt,topsep=4pt]
    \item Airflow initiates the relevant workflow from the local control plane;
    \item local producers read the source datasets and publish records to Kafka;
    \item topic partitions are hosted across the brokers on VM1, VM2, and VM3;
    \item bronze consumers read the topics and write source-specific \texttt{.jsonl} batches to HDFS;
    \item the NameNode on VM2 coordinates block persistence across the three DataNodes;
    \item Spark jobs enter through the remote submission path on VM2, with executors allocated across registered workers;
    \item Spark reads bronze data, writes raw Iceberg tables, and registers their metadata through Hive Metastore on VM2;
    \item the embedded Great Expectations tasks validate the raw lakehouse tables and persist validation results and Data Docs on the control plane;
    \item dbt-spark is initiated from the control plane while its transformation work executes through the Spark cluster;
    \item the Trino coordinator on VM2 plans analytical queries whose fragments are executed by workers on VM1 and VM3;
    \item Prometheus collects distributed service metrics and Grafana exposes centralized monitoring views.
\end{enumerate}

The resulting path crosses explicit machine and service boundaries for transport, storage, processing, metadata governance, querying, transformation, validation, and supervision. The platform therefore does not collapse its data path into a single-machine ETL process.

\begin{figure}[H]
    \centering
    \begin{tikzpicture}[node distance=0.72cm,font=\small]
  \node[ado box={12.80cm}{adograyfill}, minimum height=1.25cm, inner sep=5pt] (control)
    {\textbf{\normalsize Local control plane}\\[-1pt]
     \footnotesize Local workstation: dataset preparation and Airflow orchestration trigger};

  \node[ado box={12.80cm}{adotealfill}, minimum height=1.25cm, inner sep=5pt, below=of control] (kafka)
    {\textbf{\normalsize Kafka transport cluster}\\[-1pt]
     \footnotesize \textbf{VM1:} broker 1 \quad $\vert$ \quad
     \textbf{VM2:} broker 2 \quad $\vert$ \quad
     \textbf{VM3:} broker 3};

  \node[ado box={12.80cm}{adobluefill}, minimum height=1.25cm, inner sep=5pt, below=of kafka] (hdfs)
    {\textbf{\normalsize HDFS bronze persistence}\\[-1pt]
     \footnotesize \textbf{VM2:} NameNode + DataNode 2 \quad $\vert$ \quad
     \textbf{VM1/VM3:} DataNodes 1 and 3};

  \node[ado box={12.80cm}{adoorangefill}, minimum height=1.25cm, inner sep=5pt, below=of hdfs] (spark)
    {\textbf{\normalsize Distributed Spark compute}\\[-1pt]
     \footnotesize \textbf{VM2:} master + submit client \quad $\vert$ \quad
     \textbf{VM1/VM2/VM3:} workers 1--3};

  \node[ado box={12.80cm}{adopurplefill}, minimum height=1.55cm, inner sep=5pt, below=of spark] (lakehouse)
    {\textbf{\normalsize Governed lakehouse and SQL exposure}\\[-1pt]
     \footnotesize Iceberg tables on HDFS \quad $\vert$ \quad
     \textbf{VM2:} Hive Metastore + Trino coordinator\\[-1pt]
     \footnotesize \textbf{VM1/VM3:} Trino workers};

  \draw[ado arrow, draw=black!60, line width=1.05pt]
    (control) -- node[ado annotation, font=\scriptsize, text=black!60, fill=white, inner sep=1pt, right=2pt] {publish} (kafka);
  \draw[ado arrow, draw=black!60, line width=1.05pt]
    (kafka) -- node[ado annotation, font=\scriptsize, text=black!60, fill=white, inner sep=1pt, right=2pt] {persist} (hdfs);
  \draw[ado arrow, draw=black!60, line width=1.05pt]
    (hdfs) -- node[ado annotation, font=\scriptsize, text=black!60, fill=white, inner sep=1pt, right=2pt] {submit and execute} (spark);
  \draw[ado arrow, draw=black!60, line width=1.05pt]
    (spark) -- node[ado annotation, font=\scriptsize, text=black!60, fill=white, inner sep=1pt, right=2pt] {write, register, query} (lakehouse);
\end{tikzpicture}

    \caption{Machine-level flow across the local control and three-node data planes.}
    \label{fig:chapter6-machine-data-flow}
\end{figure}

\section{Distributed Spark Execution}
Spark provides the principal distributed-processing mechanism of the platform. The Spark master is hosted on VM2, while executor capacity is registered through workers on VM1, VM2, and VM3. Processing capacity is therefore not bound to the coordination node.

The Spark master interface exposes the registered worker pool, and applications allocate executors according to workload requirements and available capacity. Processing tasks can consequently be scheduled beyond the coordinator node itself; the corresponding runtime evidence is examined in Chapter VIII.

The deployment also includes a Spark submit client, Spark Thrift Server, and Spark History Server on VM2. These components preserve the separation between cluster entry, SQL access, execution, and historical inspection \cite{apache_spark_docs}:
\begin{itemize}
    \item Spark jobs are launched from the control plane through an SSH-based remote submission path to VM2;
    \item transformation-oriented SQL access targets the distributed Spark environment through the Thrift interface;
    \item application history is retained through the History Server for later inspection;
    \item execution is coordinated centrally while tasks are carried out by workload-selected workers.
\end{itemize}

Executor allocation is workload-dependent. A registered worker may receive no tasks during a light application, because distributed execution does not imply equal simultaneous utilization of every node. The relevant criterion is that the cluster can allocate work beyond VM2 when the workload and available resources require it.

\begin{figure}[H]
    \centering
    \begin{tikzpicture}[x=1cm,y=1cm,font=\small]
  \node[ado box={5.80cm}{adotealfill}, minimum height=1.25cm, inner sep=5pt] (master) at (0,0)
    {\textbf{\normalsize Spark master --- VM2}\\[-1pt]
     \footnotesize worker registration and task scheduling};

  \node[ado group, minimum width=14.20cm, minimum height=2.85cm, anchor=north] (pool) at (0,-1.35) {};
  \node[font=\bfseries\small, anchor=west] at ($(pool.north west)+(0.32,-0.35)$)
    {Registered Spark worker pool};
  \node[font=\scriptsize, text=adotextsoft, anchor=east] at ($(pool.north east)+(-0.32,-0.35)$)
    {active allocation depends on workload};

  \node[ado box={3.45cm}{adoorangefill}, minimum height=1.25cm, inner sep=5pt] (w1) at (-4.45,-3.10)
    {\textbf{\small Worker 1 --- VM1}\\[2pt]
     \footnotesize executor capacity};
  \node[ado box={3.45cm}{adoorangefill}, minimum height=1.25cm, inner sep=5pt] (w2) at (0,-3.10)
    {\textbf{\small Worker 2 --- VM2}\\[2pt]
     \footnotesize executor capacity};
  \node[ado box={3.45cm}{adoorangefill}, minimum height=1.25cm, inner sep=5pt] (w3) at (4.45,-3.10)
    {\textbf{\small Worker 3 --- VM3}\\[2pt]
     \footnotesize executor capacity};

  \coordinate (dispatch) at (0,-2.20);
  \draw[ado arrow, draw=black!60, line width=1.05pt]
    (master.south) -- node[font=\scriptsize, text=black!60, fill=white, inner sep=1pt, right=2pt] {schedule} (dispatch);
  \draw[draw=black!55, line width=0.9pt] (-4.45,-2.20) -- (4.45,-2.20);
  \draw[ado arrow, draw=black!60, line width=1.05pt] (-4.45,-2.20) -- (w1.north);
  \draw[ado arrow, draw=black!60, line width=1.05pt] (0,-2.20) -- (w2.north);
  \draw[ado arrow, draw=black!60, line width=1.05pt] (4.45,-2.20) -- (w3.north);

  \node[ado box={12.80cm}{adopurplefill}, minimum height=1.25cm, inner sep=5pt] (result) at (0,-4.95)
    {\textbf{\normalsize Distributed application execution}\\[-1pt]
     \footnotesize tasks run on workload-selected workers, including capacity beyond VM2};

  \draw[ado dashed arrow, draw=black!55, line width=0.9pt]
    (w1.south) -- ($(result.north)+(-4.20,0)$);
  \draw[ado dashed arrow, draw=black!55, line width=0.9pt]
    (w2.south) -- (result.north);
  \draw[ado dashed arrow, draw=black!55, line width=0.9pt]
    (w3.south) -- ($(result.north)+(4.20,0)$);
\end{tikzpicture}

    \caption{Workload-dependent Spark execution across VM1, VM2, and VM3.}
    \label{fig:chapter6-distributed-spark-execution}
\end{figure}

\section{Operational Packaging and Deployment Bundles}
The deployment is supported by explicit operational packaging rather than undocumented manual steps. Node-specific bundles, startup logic, reset logic, and preflight checks capture the operational dependencies of the platform and support reproducible execution \cite{docker_docs}.

\begin{table}[!htbp]
    \centering
    \footnotesize
    \renewcommand{\arraystretch}{1.16}
    \arrayrulecolor{tablelineblue}
    \rowcolors{2}{tablebodyblue}{white}
    \begin{tabularx}{\textwidth}{|Y|Y|Y|}
        \hline
        \rowcolor{tableheaderblue}
        \textcolor{white}{\textbf{Deployment artifact type}} & \textcolor{white}{\textbf{Role in operations}} & \textcolor{white}{\textbf{Operational benefit}} \\
        \hline
        Bundle management scripts & prepare and synchronize node-specific deployment content & maintain consistent configuration across repeated deployments \\
        \hline
        Preflight checks & verify ports, dependencies, peer endpoints, and critical services & identify blocking conditions before service startup \\
        \hline
        Reset scripts & clean previous runtime state and support reproducible reruns & restore a controlled baseline for validation cycles \\
        \hline
        Startup scripts & launch services in the required sequence & formalize the operational lifecycle of the platform \\
        \hline
        Access-point documentation & record URLs, ports, and secured entry paths & preserves the traceability of operational interfaces \\
        \hline
    \end{tabularx}
    \caption{Operational deployment artifacts used in the project}
    \label{tab:chapter6-deployment-artifacts}
\end{table}

The deployment bundles translate the architectural service map into machine-specific operational units. Each node receives the configuration, compose definitions, mounted service files, and startup logic required for its assigned role. Keeping these elements together reduces configuration drift between repeated deployments and makes the relationship between a service, its node, and its exposed interfaces explicit. Bundle preparation is therefore a reproducibility mechanism rather than a simple file-transfer step.

Preflight verification acts as a gate before any distributed workflow is launched. It checks the conditions on which later services depend, including required ports, peer reachability, expected endpoints, deployment content, and critical prerequisites. Detecting these conditions before startup avoids partial cluster states in which a worker is running but cannot reach its coordinator, or an orchestration task begins while the storage, catalog, or execution layer is still unavailable. The checks also provide a clearer failure boundary because deployment problems can be separated from data-pipeline failures.

The startup order is intentional. VM2 is started first because it hosts the coordination-heavy services: the Kerberos authority, HDFS NameNode, Hive Metastore, Spark master, and Trino coordinator establish identities, namespaces, metadata, scheduling, and query coordination for the rest of the platform. VM1 and VM3 are then started to register their storage, transport, compute, and query-worker capacity against those shared services. Finally, the local control plane is activated only after data-plane readiness has been confirmed, preventing Airflow and the monitoring stack from interpreting normal initialization delays as operational failures. This sequence follows service dependencies rather than an arbitrary launch order.

\clearpage

\begin{figure}[H]
    \centering
    \begin{tikzpicture}[node distance=0.55cm]
    \node[ado box=adograyfill, text width=10.2cm] (bundle)
    {\adoboxcontent{Bundle preparation}{node-specific deployment content + synchronized configuration}};
    \node[ado box=adobluefill, text width=10.2cm, below=of bundle] (preflight)
    {\adoboxcontent{Preflight verification}{ports + peer reachability + dependency checks}};
    \node[ado box=adotealfill, text width=10.2cm, below=of preflight] (vm2)
    {\adoboxcontent{Priority startup of VM2}{coordination-heavy services launched first}};
    \node[ado box=adoorangefill, text width=10.2cm, below=of vm2] (workers)
    {\adoboxcontent{Cluster extension}{VM1 and VM3 add storage, broker, and worker capacity}};
    \node[ado box=adopurplefill, text width=10.2cm, below=of workers] (cp)
    {\adoboxcontent{Control-plane activation}{Airflow + monitoring + access services start after cluster stability}};
    \draw[ado arrow] (bundle) -- (preflight);
    \draw[ado arrow] (preflight) -- (vm2);
    \draw[ado arrow] (vm2) -- (workers);
    \draw[ado arrow] (workers) -- (cp);
\end{tikzpicture}

    \caption{Dependency-aware deployment lifecycle and startup order.}
    \label{fig:chapter6-deployment-lifecycle}
\end{figure}
\FloatBarrier

\section{Conclusion}
This chapter clarified that the deployment of \projectname{} is intentionally distributed at both the service level and the machine level. The local control plane supervises the environment without collapsing it into a single node, VM2 plays the role of coordination leader, and VM1 and VM3 extend storage, compute, and worker capacity. Data movement, processing, and analytical access therefore take place across explicit multi-node boundaries.

The deployment is reproducible through node-specific bundles, reset logic, startup scripts, and preflight checks. It therefore constitutes an implemented operational topology rather than only a conceptual design. Chapter VII examines the security controls that harden this deployment and protect its critical access paths.
